\documentclass[12pt]{article}
\usepackage[utf8]{inputenc}
\usepackage[T1]{fontenc}
\usepackage{hyperref}
\title{Kinnauri, Chhoyul Language Resources for K-12 Curriculum Development}
\author{Compiled by Claude Sonnet 5 (claude-sonnet-5) for Andr\'as Kornai}
\date{}
\begin{document}
\maketitle

\section{Introduction}
Chhoyul Kinnauri (ISO 639-3: tpq) is a variety of Kinnauri spoken in Himachal Pradesh, India. It is Sino-Tibetan and endangered. Sources give different speaker counts: 610 (Global Recordings Network) and 800 (Joshua Project). Alternate names listed include Nesang, Nyamkad and Tukpa. Search found very little specific to this variety. Do not confuse it with Chhitkul-R\=akchham (cik), which is a separate language.

\section{Primary Educational Resources}
No primers, textbooks or school materials for Chhoyul were found.

\section{Linguistic and Reference Materials}
\begin{itemize}
\item Glottolog notes Kinnauri as definitely endangered; the Kuno-Charang dialect page is: \url{https://glottolog.org/resource/languoid/id/tukp1239}
\item The Linguistic Landscape of the Indian Himalayas (open-access book; I did not check whether it covers Chhoyul): \url{https://library.oapen.org/handle/20.500.12657/76653}
\item Endangered Languages Project bibliography for a related entry: \url{https://www.endangeredlanguages.com/lang/4373/bibliography}
\end{itemize}

\section{Religious and Community Resources}
\begin{itemize}
\item Joshua Project profile of the Kanaura people group: \url{https://joshuaproject.net/people_groups/12487/IN}
\end{itemize}

\section{Digital and Audio Resources}
\begin{itemize}
\item Global Recordings Network page for tpq (evangelistic audio recordings): \url{https://globalrecordings.net/en/language/tpq}
\end{itemize}
The Chhitkul-R\=akchham collection at UNT is a different language and is not a Chhoyul resource: \url{https://digital.library.unt.edu/explore/collections/CRLR/}

\section{Curriculum Development Context}
Very little is documented for Chhoyul. Any curriculum would need primary documentation with speakers first, and could draw on materials for standard Kinnauri.
\end{document}